\documentclass[11pt]{article}
\usepackage[margin=1.1in]{geometry}
\usepackage{amsmath,amssymb}
\usepackage{booktabs}
\usepackage{array}
\usepackage{xcolor}
\usepackage{hyperref}
\usepackage{microtype}
\usepackage{parskip}
\usepackage{titlesec}
\usepackage{enumitem}
\usepackage{fancyhdr}
\usepackage{graphicx}

\definecolor{dkred}{RGB}{180,30,30}
\definecolor{dkblue}{RGB}{20,60,140}
\definecolor{dkgreen}{RGB}{20,110,60}
\definecolor{muted}{RGB}{80,80,95}

\hypersetup{
  colorlinks=true,
  linkcolor=dkblue,
  urlcolor=dkblue,
  citecolor=dkblue
}

\titleformat{\section}{\large\bfseries\color{dkblue}}{}{0em}{}
\titleformat{\subsection}{\normalsize\bfseries}{}{0em}{}

\pagestyle{fancy}
\fancyhf{}
\rhead{\small\color{muted}MIRADOR Clinical Brief · September 2026}
\lhead{\small\color{muted}The Keske Method — Pediatric AHO}
\rfoot{\small\color{muted}\thepage}
\renewcommand{\headrulewidth}{0pt}

% ─────────────────────────────────────────────────────────────────────────────
\begin{document}

\begin{center}
  {\LARGE\bfseries The Keske Method}\\[6pt]
  {\large A Geometric Framework for Antibiotic Selection\\
  in Pediatric MRSA Acute Hematogenous Osteomyelitis}\\[10pt]
  {\color{muted}\small
  Bee Rosa Davis, MS (Digital Forensics, Brown University) · Applied Mathematics ·
  Independent Researcher \\
  ORCID: 0009-0009-8034-4308 · \href{mailto:bee_davis@alumni.brown.edu}{bee\_davis@alumni.brown.edu} ·
  \href{https://doi.org/10.5281/zenodo.18511755}{doi:10.5281/zenodo.18511755}}\\[4pt]
  {\color{dkred}\small\itshape
  Named for Steven Keske — 6 years, 5 antibiotics, 4 surgeries, still fighting.}
\end{center}

\vspace{6pt}

% ─────────────────────────────────────────────────────────────────────────────
\section{The Clinical Gap}

Acute hematogenous osteomyelitis (AHO) caused by community-associated MRSA
(USA300 pulsotype) remains one of the most clinically treacherous pediatric
infections. The United States sees approximately 70,000 pediatric musculoskeletal
infections annually, with MRSA-AHO driving the most severe complications:
septic shock, toxic shock syndrome (TSS), septic pulmonary emboli, pathologic
fractures, and chronic relapsing disease requiring multi-year antibiotic courses.

The literature documents a reproducible failure mode: patients cycle through
five or more antibiotic regimens over months to years, each serologically
adequate, each clinically insufficient. Treatment guidelines remain empirical,
and chronic biofilm-associated MRSA-AHO is widely regarded as refractory to
standard antimicrobial therapy.

\textbf{The Keske Method proposes that this failure is not primarily
microbiological — it is geometric. The drugs exist. The standard pharmacokinetic
models used to deploy them were built for blood, not bone. The math to deploy
them correctly in deep-tissue sanctuaries is not part of routine clinical
practice.}

% ─────────────────────────────────────────────────────────────────────────────
\section{Governing Equation}

The central question in antibiotic selection is: \textit{will this drug reach a
concentration high enough at the site of infection to kill the bacteria there?}
The Keske Method formalizes that question as a single score, $C_{\text{bone}}$,
computed from every known barrier between the IV bag and the bacteria inside
bone. Higher $C$ means fewer barriers, better penetration, higher probability of
cure. The working cure threshold $C_{\text{bone}} \geq 5$ corresponds to
achieving a bone concentration at or above the published MBEC for the dominant
pathogen — the concentration required to eradicate biofilm, not merely suppress
planktonic growth (Zimmerli et al., NEJM 2004; Stewart, Microbiol Mol Biol Rev
2015). This threshold is a modeling anchor derived from published data; it is
not itself prospectively validated (see \S\ref{sec:validation}).

The Keske Method derives from the MIRADOR framework (Manifold-Informed Rational
Architecture for Drug-Organism Response). The governing equation is:
\[
  C = \frac{\tau}{K}
\]
$\tau$ (pharmacophoric potential) captures the drug's intrinsic activity: its
binding affinity to the target (PBP2a for MRSA), aqueous solubility, oral or IV
bioavailability, and protein binding. Clinically, $\tau$ is proportional to the
drug's theoretical AUC$_{0-24}$/MIC under ideal distribution conditions — the
ceiling of what the drug can achieve if all barriers are removed.

$K$ (total curvature) is the sum of all impedance factors that reduce the drug's
effective concentration at the infection site. Clinically, $K$ is what separates
adequate serum levels from adequate bone levels. A drug with $K = 1$ loses 50\%
of its activity to barriers; a drug with $K = 9$ loses 90\%.

For standard bacteremia, $K = K_{\text{ADMET}}$ (absorption, distribution,
metabolism, excretion, toxicity). For pediatric AHO, three additional terms
dominate:
\[
  \boxed{K_{\text{bone}} = K_{\text{ADMET}} + K_{\text{penetration}} + K_{\text{biofilm}} + K_{\text{reservoir}}}
\]

Each term is clinically grounded, computationally specified, and sourced from
published pharmacokinetic and microbiological literature. Every number in this
brief is produced by the open-source MIRADOR engine and is reproducible from the
inputs shown (\S\ref{sec:whatis}).

% ─────────────────────────────────────────────────────────────────────────────
\section{The Four Curvature Layers}

\subsection{Layer 1: Pediatric Patient Manifold}

Children are not small adults. The Keske Method replaces adult pharmacokinetic
assumptions with pediatric-specific models:

\begin{align}
  \text{GFR}_{\text{ped}} &= \frac{0.413 \times h_{\text{cm}}}{S_{\text{cr}}} \quad
    \text{(Schwartz equation)} \\[4pt]
  \text{CL}_{\text{ped}} &= \text{CL}_{\text{adult}} \cdot
    \left(\frac{W}{70}\right)^{0.75} \quad
    \text{(allometric scaling)} \\[4pt]
  V_d &= V_{d,\text{adult}} \cdot \left(\frac{W}{70}\right)^{1.0}
    \cdot \bigl[1 + 0.003 \cdot \max(\text{CRP} - 100, 0)\bigr]
\end{align}

The CRP-dependent $V_d$ expansion (20--40\% at CRP $>100$ mg/L) reflects
capillary leak in systemic inflammation (Graziani et al., JAC 1988).

\subsection{Layer 2: Bone Penetration Curvature}

The critical barrier that standard dosing models ignore is the serum-to-bone
partitioning ratio $R_{\text{bone}}$. Published bone:serum ratios are
drug-specific and non-trivial:

\begin{table}[h!]
\centering
\small
\begin{tabular}{lccl}
\toprule
\textbf{Drug} & \textbf{Bone:Serum} & \textbf{$K_{\text{penetration}}$} & \textbf{Source} \\
\midrule
Vancomycin   & 0.10--0.30 & 2.33--9.00 & Graziani et al., JAC 1988 \\
Ceftaroline  & 0.20--0.40 & 1.50--4.00 & Riccobene et al., AAC 2014 \\
Linezolid    & 0.40--0.60 & 0.67--1.50 & Rana et al., JAC 2002 \\
Clindamycin  & 0.30--0.75 & 0.33--2.33 & Feigin et al., Pediatr Infect 1995 \\
Rifampin     & 0.20--0.50 & 1.00--4.00 & Currier et al., JBJS 1979 \\
Daptomycin   & 0.10--0.20 & 4.00--9.00 & Traunmuller et al., AAC 2010 \\
\bottomrule
\end{tabular}
\caption{Bone penetration ratios and resulting curvature contributions.
$K_{\text{penetration}} = (1/R_{\text{bone}}) - 1$.}
\end{table}

The penetration curvature is:
\[
  K_{\text{penetration}} = \frac{1}{R_{\text{bone,eff}}} - 1
\]

Inflamed bone is more vascular; CRP $> 100$ mg/L increases $R_{\text{bone}}$:
\[
  R_{\text{bone,eff}} = R_{\text{bone}} \cdot
    \bigl[1 + 0.006 \cdot \max(\text{CRP} - 100,\, 0)\bigr], \quad
    \text{capped at } 2 \times R_{\text{bone}}
\]

This is why vancomycin attains adequate serum levels — AUC/MIC $>400$ in many
patients — and still fails in chronic osteomyelitis: the serum target and the
bone target are not the same target.

\subsection{Layer 3: Biofilm Resistance Curvature}

MRSA in established osteomyelitis forms biofilm — structured communities of
phenotypically dormant bacteria embedded in extracellular matrix. Biofilm MICs
are 100- to 1000-fold higher than planktonic MICs (Stewart, 2015). The minimum
biofilm eradication concentration (MBEC) determines effective drug activity:

\[
  K_{\text{biofilm}} = \log_{10}\!\left(\frac{\text{MBEC}}{\text{MIC}}\right)
\]

\begin{table}[h!]
\centering
\small
\begin{tabular}{lcccl}
\toprule
\textbf{Drug} & \textbf{MIC ($\mu$g/mL)} & \textbf{MBEC ($\mu$g/mL)} & \textbf{$K_{\text{biofilm}}$} & \textbf{Source} \\
\midrule
Vancomycin   & 1.0   & 512  & 2.71 & Parra-Ruiz et al., 2012 \\
Clindamycin  & 0.25  & 64   & 2.41 & LaPlante \& Rybak, 2004 \\
Ceftaroline  & 1.0   & 128  & 2.11 & Barber et al., AAC 2015 \\
Linezolid    & 2.0   & 256  & 2.11 & Parra-Ruiz et al., 2012 \\
Daptomycin   & 0.5   & 32   & 1.81 & Stewart, 2015 \\
Rifampin     & 0.008 & 0.5  & 1.80 & Zimmerli et al., NEJM 2004 \\
\bottomrule
\end{tabular}
\caption{MIC, MBEC, and biofilm curvature by drug. Rifampin has the lowest
$K_{\text{biofilm}}$ of any anti-staphylococcal agent — the geometric basis
for its role in combination therapy for chronic AHO.}
\end{table}

Biofilm probability is inferred from the clinical timeline:
\[
  p_{\text{biofilm}} =
  \begin{cases}
    0.20 & \text{acute} \;(<\!2\;\text{weeks}) \\
    0.60 & \text{subacute} \;(2\;\text{weeks–}3\;\text{months}) \\
    0.95 & \text{chronic} \;(>\!3\;\text{months})
  \end{cases}
  \qquad
  K_{\text{biofilm,eff}} = p_{\text{biofilm}} \cdot K_{\text{biofilm}}
\]

The MIC reported on an antibiogram is the planktonic MIC. The bacteria causing
chronic osteomyelitis are in biofilm, with an effective MIC 100- to 1000-fold
higher. The Keske Method computes this from published MBEC data and weights it
by the probability that biofilm is present.

\subsection{Layer 4: Multi-Reservoir Geometry}

MRSA-AHO distributes across three anatomically distinct reservoirs with
different antibiotic accessibility:

\begin{table}[h!]
\centering
\small
\begin{tabular}{p{3.2cm} p{2.8cm} p{2.8cm} p{3.5cm}}
\toprule
\textbf{Reservoir} & \textbf{Location} & \textbf{Accessibility} & \textbf{Intervention} \\
\midrule
Soft tissue abscess & Periosteal / subperiosteal & High (vascularized) & Surgical drainage \\
Bone matrix biofilm & Cortical / cancellous surfaces & Moderate (penetration-limited) & Debridement \\
Intracellular SCVs & Inside osteocytes & Very low & Rifampin; no surgical access \\
\bottomrule
\end{tabular}
\caption{Three-reservoir model. Small colony variants (SCVs) inside osteocytes
drive relapse in chronic AHO. They are not reached by surgery and are accessible
only to drugs with intracellular penetration — principally rifampin.}
\end{table}

The reservoir curvature term, per drug:
\[
  K_{\text{reservoir}} = (1 - D_{\text{drain}}) \cdot w_{\text{SAC}}
    + (1 - D_{\text{debride}}) \cdot K_{\text{penetration}}
    + 0.8 \cdot f_{\text{intra}} \cdot m_{\text{rif}}
\]
where $D_{\text{drain}}$, $D_{\text{debride}}$ are surgical clearance fractions
(0–1), $f_{\text{intra}}$ is the intracellular bacterial fraction,
$w_{\text{SAC}} = 0.5$, and $m_{\text{rif}} = 0.4$ when rifampin is present in
the regimen (reflecting its documented osteoblast penetration; Zimmerli 2004),
else $m_{\text{rif}} = 1.0$.

% ─────────────────────────────────────────────────────────────────────────────
\section{Applied to a Clinical Scenario}

\textbf{Patient:} 10-year-old, 32 kg. MRSA-AHO of the hip, femur, and knee.
Sepsis and TSS on presentation; chronic recurrent course. Multiple antibiotics
trialed without durable remission.

\textbf{Input parameters:}
\[
  \text{CRP} = 250 \;\text{mg/L}, \quad
  p_{\text{biofilm}} = 0.95, \quad
  f_{\text{intra}} = 0.40, \quad
  D_{\text{drain}} = 0.90, \quad
  D_{\text{debride}} = 0.80
\]

\textbf{Computed $C_{\text{bone}}$ for each drug — monotherapy.} Every column is
produced by the engine; $K_{\text{total}} = K_{\text{ADMET}} + K_{\text{pen}} +
K_{\text{bio}} + K_{\text{res}}$ and $C_{\text{bone}} = \tau / K_{\text{total}}$.

\begin{table}[h!]
\centering
\small
\begin{tabular}{lccccccc}
\toprule
\textbf{Drug} & \textbf{$\tau$} & \textbf{$K_{\text{ADMET}}$} & \textbf{$K_{\text{pen}}$} & \textbf{$K_{\text{bio}}$} & \textbf{$K_{\text{res}}$} & \textbf{$K_{\text{total}}$} & \textbf{$C_{\text{bone}}$} \\
\midrule
Linezolid    & 12 & 0.40 & 0.05 & 2.00 & 0.38 & 2.84 & 4.23 \\
Daptomycin   & 24 & 0.60 & 2.51 & 1.72 & 0.87 & 5.70 & 4.21 \\
Ceftaroline  & 12 & 0.67 & 0.75 & 2.00 & 0.52 & 3.95 & 3.04 \\
Rifampin$^\dagger$ &  8 & 0.50 & 0.50 & 1.71 & 0.28 & 2.99 & 2.68 \\
Clindamycin  &  8 & 0.50 & 0.00 & 2.29 & 0.37 & 3.16 & 2.53 \\
Vancomycin   & 12 & 0.50 & 1.63 & 2.57 & 0.70 & 5.40 & 2.22 \\
\bottomrule
\end{tabular}
\caption{Cure threshold: $C_{\text{bone}} \geq 5.0$. \textbf{No monotherapy
achieves it} in chronic MRSA-AHO with confirmed biofilm — the best single agent
reaches only $4.23$. This matches the documented clinical course for chronic AHO,
in which serologically adequate monotherapy fails. $^\dagger$Rifampin is never
used as monotherapy (rpoB resistance emerges in $\sim$80\% of cases); its value
is shown only as a combination component.}
\end{table}

\subsection*{Combination therapy}

No single agent reaches threshold, so the method scores a two-drug regimen. Each
drug is first reduced to its own bone coherence $C = \tau / K_{\text{pathway}}$;
the two are then combined by their \emph{measured pharmacodynamic interaction} —
not by a fixed multiplier that always helps. Three interaction types are
representable, so an unfavorable pairing is not silently rewarded:
\[
  C_{\text{combo}} =
  \begin{cases}
    C_A + C_B & \text{additivity (Bliss independence)} \\
    \max(C_A, C_B) + \delta & \text{synergy} \\
    \min(C_A, C_B) \cdot (1 - \delta) & \text{antagonism}
  \end{cases}
\]
The interaction type and its magnitude $\delta$ are \emph{evidence inputs}
(checkerboard / FICI / Bliss), supplied per isolate rather than assumed.

For \textbf{ceftaroline + rifampin} at this geometry — with rifampin's
intracellular reservoir term reduced ($m_{\text{rif}} = 0.4$) — the per-drug bone
coherences are $C_{\text{ceftaroline}} = 3.20$ and $C_{\text{rifampin}} = 2.68$.
Under Bliss independence:
\[
  C_{\text{bone}}(\text{cef} + \text{rif}) = 3.20 + 2.68 =
  {\color{dkgreen}\mathbf{5.87}} \;\geq\; 5.
\]

\textbf{This is the first regimen in this patient's geometry to cross the cure
threshold} — where vancomycin, the standard-of-care agent, scores $2.22$. The
result is a $2.6\times$ improvement in therapeutic coherence, and it comes from
the two barriers no single agent overcomes: bone penetration and rifampin's
access to the intracellular reservoir.

The result is explicitly \emph{contingent on the interaction}: additivity or
synergy places the regimen at or above threshold, while strain-level antagonism
— which has been reported for rifampin combinations in some isolates (Barber et
al., AAC 2015) — would place it below. Determining the interaction for a given
isolate, by checkerboard or FICI, is therefore the decisive measurement. For
reference, an earlier ``parallel-conductance'' combination heuristic is retired
in favor of this evidence-typed model, which cannot manufacture benefit from an
inactive or antagonistic partner.

% ─────────────────────────────────────────────────────────────────────────────
\section{Clinical Implication: What the Math Tells Us}

\begin{enumerate}[leftmargin=1.4em]
  \item \textbf{Serum adequacy is not bone adequacy.} Vancomycin achieving
        AUC/MIC $> 400$ in serum is the wrong target when $R_{\text{bone}}$ is
        low and $p_{\text{biofilm}} = 0.95$. That target was derived from
        bacteremia data, not osteomyelitis data. The Keske Method quantifies the
        bone-specific deficit prospectively rather than inferring it
        retrospectively from treatment failure.

  \item \textbf{Biofilm is the dominant curvature term in chronic disease, and
        antibiograms do not measure it.} At $p_{\text{biofilm}} \geq 0.8$, no
        monotherapy crosses the cure threshold; a combination is what brings a
        regimen across it. Rifampin's low $K_{\text{biofilm}}$ and unique
        intracellular access make it the natural partner — consistent with its
        established role in refractory staphylococcal bone and implant infection
        (Zimmerli 2004). Whether a specific pairing reaches threshold depends on
        the measured drug-drug interaction, which the model takes as an input
        rather than assuming.

  \item \textbf{Surgical intervention has a quantifiable pharmacodynamic effect.}
        Drainage and debridement reduce $K_{\text{reservoir}}$ and secondarily
        raise $R_{\text{bone,eff}}$ by reducing inflammatory mass. The model
        computes the specific $D_{\text{drain}}$ and $D_{\text{debride}}$ values
        that, in combination with a given drug pair, bring $C_{\text{bone}}$ to
        $\geq 5$ — an operable, patient-specific surgical target rather than a
        generic recommendation.

  \item \textbf{Intracellular reservoirs are the remission barrier.} At high
        $f_{\text{intra}}$, a regimen cannot hold threshold without a drug that
        has documented intracellular penetration; rifampin and linezolid are the
        principal anti-MRSA agents in that class. Sequential design — a
        cell-wall-active agent first to reduce biofilm mass, then a rifampin
        combination to address SCVs — is the direction the geometry supports.
\end{enumerate}

% ─────────────────────────────────────────────────────────────────────────────
\section{What This Is — and What It Is Not}\label{sec:whatis}

The Keske Method is an \textbf{investigational mathematical framework}. It is not
an FDA-cleared diagnostic, dosing calculator, or prescribing tool. It is
presented strictly for theoretical evaluation, retrospective study, and
hypothesis generation by qualified clinical researchers. It does not replace
susceptibility testing, infectious disease consultation, or clinical judgment. It
takes no action autonomously and makes no therapeutic recommendation.

What it provides: a systematic, patient-specific, mathematically grounded answer
to \textit{``why didn't that antibiotic work?''} — and a directional answer to
\textit{``what geometry do we need to reach threshold?''} Every figure in this
brief is computed by the open engine and is reproducible from the stated inputs.

The live interactive tool (MIRADOR — Keske Edition) is publicly available at:\\
\url{https://mirador.vercel.app/\#keske}\\[2pt]
Repository and validation results: \href{https://doi.org/10.5281/zenodo.18511755}{doi:10.5281/zenodo.18511755}

\vspace{6pt}
\textbf{Medical and clinical patent applications} (selected; Davis, B.R., sole
inventor; serial numbers as filed):

\begin{enumerate}[leftmargin=1.6em, itemsep=2pt]
  \item Davis, B.R.\ \textit{MIRADOR: Manifold-Informed Rational Architecture for
        Drug-Organism Response — Geometric Optimization of Therapeutic Coherence
        with Resistance Prediction via Escape Geodesics on Drug-Target Fiber
        Bundles.} U.S.\ non-provisional application (serial no.\ on file with the
        USPTO), 2026; claiming priority to a related provisional.

  \item Davis, B.R.\ \textit{TESSERA: Transfer-Enabled Segmented Surveillance of
        Emergent Resistance in Antimicrobials.}
        U.S.\ Provisional Patent Application No.~63/933,236.

  \item Davis, B.R.\ \textit{Genomic Evolution on Davis-Embedded Surfaces for
        Integrated Cancer Detection.}
        U.S.\ Provisional Patent Application No.~63/931,845.

  \item Davis, B.R.\ \textit{HERALD: High-Resolution Early Recognition of
        Antigenic Landscape Divergence.}
        U.S.\ Provisional Patent Application No.~63/919,595.

  \item Davis, B.R.\ \textit{Geometric Phase Transition Detection and
        Curvature-Based Resistance Emergence Threshold Identification.}
        U.S.\ Provisional Patent Application No.~63/926,936.
\end{enumerate}

% ─────────────────────────────────────────────────────────────────────────────
\section{Proposed Retrospective Validation}\label{sec:validation}

The $C_{\text{bone}} \geq 5$ threshold and the curvature weights used throughout
this model are derived from published pharmacokinetic, microbiological, and
clinical data. They have not been prospectively validated in a clinical cohort.
We propose the following validation design for researchers with access to
institutional records:

\textbf{Study design:} Multi-center retrospective cohort of $\geq 100$ cases of
pediatric MRSA-AHO (age $< 18$, culture-confirmed MRSA, minimum 12-month
follow-up), classified as treatment success (remission at 12 months) or failure
(relapse, recurrence, or chronic suppression).

\textbf{Primary analysis:} For each historical case, input the available clinical
variables (admission CRP, infection duration, surgical interventions, antibiotic
sequence) into the MIRADOR engine to compute $C_{\text{bone}}$ for the regimen
actually used. Test whether $C_{\text{bone}} < 5$ predicts clinical failure
better than standard metrics (vancomycin AUC/MIC; $\text{MIC} \leq 1\;\mu$g/mL
cutoff).

\textbf{Hypothesis:} Logistic regression will show the geometric impedance term
$K_{\text{pathway}}$ is a significant predictor of relapse independent of serum
drug exposure, with AUROC $> 0.80$ against clinical outcome.

\textbf{Data requirements:} Admission CRP, weight, estimated GFR, antibiotic(s)
used, surgical records, infection duration, and 12-month outcome — all available
in standard EHR exports, requiring no additional testing or prospective
enrollment. The engine accepts structured CSV input and is freely accessible at
\url{https://mirador.vercel.app/\#keske}.

% ─────────────────────────────────────────────────────────────────────────────
\section{Invitation}

If this framework is of interest for clinical validation, collaborative study, or
integration into pediatric AHO research protocols, I welcome the conversation.

\medskip
\noindent
\textbf{Bee Rosa Davis}\\
Applied Mathematician · Independent Researcher\\
\href{mailto:bee_davis@alumni.brown.edu}{bee\_davis@alumni.brown.edu}\\
ORCID: \href{https://orcid.org/0009-0009-8034-4308}{0009-0009-8034-4308}\\
LinkedIn: \href{https://www.linkedin.com/in/msbeedavis/}{msbeedavis}

\vspace{10pt}
{\small\color{muted}
The Keske Method is named for Steven Keske, who was diagnosed with sepsis,
toxic shock syndrome, and acute hematogenous osteomyelitis in April 2020 at age
10. He survived 10 days in the PICU, 4 surgeries, and 2 months on a PICC line.
He has experienced recurrent infections for six years. He is still fighting.
His mother Barbara shared his story publicly so that other families would know
what to look for. That story became this framework.
}

\end{document}
